\section{Conclusions}\label{sec:conclude}

\par Software Supply Chain is a crucial infrastructure of
our digitalized world. It is nevertheless susceptible to 
vulnerable/malicious code, compromised build infrastructures,
and targeted human developers. To this end, a wide variety of
LLM-based defense techniques have been proposed, where LLMs
are leveraged because of their capabilities of contextual
reasoning and code comprehension.
However, the sheer number of SSC attacks in recent years
calls for more effective defense methods.

\par This study also delves into LLM's employment for attacking
purposes. Through a rigorous analysis of collected 23 "LLM for SSC attack"
studies presented in RQ3 (Section~\ref{sec:rq3}), we conclude that
LLMs are capable of assisting attackers in breaching SSC
via vulnerable code injection or social engineering attack.

\paragraph*{Calls for ethical use of LLMs in software development}
Based on our literature review and analysis, we suggest the following
mitigation and counter-measures:

\begin{itemize}
  \item \textit{Software Developers} (Users Code Language Models) 
  should be aware of the LLM's
  tendency in providing seemingly correct code with subtle
  bugs, vulnerabilities, or even malicious code. Therefore,
  rigorous code-review and state-of-the-art linter are recommended.
  Nevertheless, the use of LLM in reporting software vulnerabilities
  must be followed by thorough investigation of veteran developers
  in case of inaccurate content.

 \item \textit{LLM providers} should be cautions of data poisoning
 attacks and low-quality training data, so that the safety of LLM-generated
  code can be improved. Furthermore, ethical alignment of LLMs
  should be implemented to counter jailbreaking or prompt injection
  attacks. 

  \item \textit{Package Managers} and \textit{Code Hosting Platforms}
  should enhance scrutiny and governance of published software packages, possibly
  employing state-of-the-art threat detectors in case of potential
  attacks from malicious code or severe vulnerabilities.
\end{itemize}

